% \iffalse meta-comment
%
% hit-final-resolution.dtx 是答辩决议页
%
% \fi
%
% \section{答辩决议页}
%
% \changes{v3.2a}{2026/09/29}{名单与决议正文照 iota-hit 是 \cs{defense} 的参数（\option{reviewers}、\option{chair}、\option{members}、
%   \option{secretary}、\option{resolution}）}
% \changes{v3.2a}{2026/09/27}{答辩决议页从贴扫描件改成排出来，照范例 \file{docx} 里那张 Word 表：
%   \pkg{docxlike} 的 \env{tabular}，六栏宽照 \texttt{tblGrid}、行高最小值 425 缇、框线半磅；格边距与格里
%   的字照 iota-hit 紧排（1\,pt、不加网格字距增量，七个字的单位名排得进一行）；合并格照 iota-hit（一行的标签顶端对齐，竖排与两个字的标签在合并出来的高度里居中，多字竖排
%   一字一格，两个字的中间空行）；决议区是接在表下的一格，排到本页剩下的整数个网格行，放不下就跨页
%   接着排。名单与决议正文由 \texttt{defense} 族给，某一块人数超了先向还有空行的块借；
%   \cs{resolution} 只剩带废弃提示的转发，排过就记闸门；人文社科类照人文范例原件那张表的列宽与五号字}
%
% 这一页原先只有两条路：给个 PDF 就贴进来，不给就贴一张空表的位图。位图是扫描来的，
% 缩放之后发虚，字体与正文两样，选不中也搜不到，英文论文照样出中文表。现在改成直接排。
%
% 数据走 \cs{hitsetup} 的 \texttt{defense} 族，一个人一条记录。表头文案走
% \texttt{term} 族，中英各一份，英文论文自动出英文表。
%
% \emph{照 Word 表排。} 范例里这是一张普通的 Word 表，所以这里用 \pkg{docxlike} 的表格：
% \env{tabular} 的格边距、竖线、横线、行的撑杆与格里的断行都换成 Word 的（见
% \file{src/docxlike/docxlike-table.dtx}），行高最小值走它的“指定高度”。合并格与决议区
% 这两样 \env{tabular} 没有，前者用悬挂的盒，后者在表下另画一个框，这两处照 iota-hit 的做法。
%    \begin{macrocode}
%<*final-resolution>
\bool_if:NT \g__hit_final_bool
  {
%    \end{macrocode}
%
% \subsection{尺寸}
%
% 范例 \file{docx}（\file{研究生学位论文书写范例}，“学位论文评阅人、答辩委员会名单及答辩决议”
% 那一章）里这张表的取值：
%
% \begin{longtblr}[caption={答辩决议表的取值}, label={tab:impl-resolution}]
%   {colspec={l X[1,l]}, width=\linewidth}
% \toprule
% 取值 & 在 \file{docx} 里 \\ \midrule
% 列宽 & \texttt{tblGrid}：675、851、1134、2126、1843、2091 缇，合 436\,bp \\
% 行高 & 前十二行 \texttt{trHeight} 425 缇（最小值），末行 2835 缇 \\
% 框线 & \texttt{tblBorders} 全是 \texttt{single}、\texttt{sz=4}，半磅 \\
% 对齐 & 表居中；格里的段居中，决议那一格左对齐；格没写 \texttt{vAlign}，顶端对齐 \\
% 字 & \texttt{Normal}，小四宋体；表头与各块的标签加粗 \\
% \bottomrule
% \end{longtblr}
%
% 表比版心宽 10.8\,bp，居中就两边各探出 5.4\,bp，外框的位置与 Word 一样。
%
% 人文社科类照《博士研究生学位论文书写范例（人文社科类）》原件里那张表（iota-hit 也照它）：
% 列宽 673、835、1106、2068、1787、2025 缇，合 424.7\,bp，比版心窄；格里的字跟那份的
% \texttt{Normal}，五号。行高、框线、对齐与理工一样。
%
% \emph{格边距与格里的字不跟 Word，跟 iota-hit：边距 1\,pt、字紧排。} Word 的格边距 0.19\,cm 下，
% “工作单位”那一栏留给字 81.35，“哈尔滨工业大学”七个字加上网格字距增量要 87.2，最常见的单位名就得
% 折行；纸质表本来就是紧排的。iota-hit 把边距收到 1\,pt、格里的字不加增量，那一栏留给字 90.15，七个字
% 84 排得下、八个字折（旧 \hithesis 用 \pkg{tabularray} 时也是 1\,pt）。“紧排”写成一个样式
% \option{resolution-table}：字体不对齐网格，Word 里就是不加增量。决议那一格照旧：边距 0.19\,cm、
% 字带增量，正文从版心左缘起。
%    \begin{macrocode}
\dim_const:Nn \c__hit_res_inset_dim { 1pt }
\dim_const:Nn \c__hit_res_row_dim  { 21.25bp }
\dim_const:Nn \c__hit_res_rule_dim { 0.5bp }
\dim_new:N \g__hit_res_wa_dim  \dim_new:N \g__hit_res_wb_dim
\dim_new:N \g__hit_res_wc_dim  \dim_new:N \g__hit_res_wd_dim
\dim_new:N \g__hit_res_we_dim  \dim_new:N \g__hit_res_wf_dim
\dim_new:N \g__hit_res_wab_dim \dim_new:N \g__hit_res_text_dim
\cs_new_protected:Npn \__hit_res_col:Nn #1#2
  { \dim_gset:Nn #1 { #2 - 2 \c__hit_res_inset_dim } }
\cs_new_protected:Npn \__hit_res_cols:nnnnnn #1#2#3#4#5#6
  {
    \__hit_res_col:Nn \g__hit_res_wa_dim  {#1}
    \__hit_res_col:Nn \g__hit_res_wb_dim  {#2}
    \__hit_res_col:Nn \g__hit_res_wc_dim  {#3}
    \__hit_res_col:Nn \g__hit_res_wd_dim  {#4}
    \__hit_res_col:Nn \g__hit_res_we_dim  {#5}
    \__hit_res_col:Nn \g__hit_res_wf_dim  {#6}
    \__hit_res_col:Nn \g__hit_res_wab_dim { #1 + #2 }
    \dim_gset:Nn \g__hit_res_text_dim { #1 + #2 + #3 + #4 + #5 + #6 - 2 \tabcolsep }
  }
\cs_new_protected:Nn \__hit_resolution_calc:
  {
    \bool_if:NTF \g__hit_hass_bool
      { \__hit_res_cols:nnnnnn { 33.65bp } { 41.75bp } { 55.30bp } { 103.40bp } { 89.35bp } { 101.25bp } }
      { \__hit_res_cols:nnnnnn { 33.75bp } { 42.55bp } { 56.70bp } { 106.30bp } { 92.15bp } { 104.55bp } }
  }
%    \end{macrocode}
%
% \subsection{行数}
%
% 每一块的行数取“填了几个人”与“纸质表上有几行”里大的那个：填了就按人数排，
% 一条不填排出来仍是与纸质表一样的空表。
%
% \textbf{某一块超了先借，借不到才占决议区的地方。} 纸质表四块的行数是 2、1、6、1，
% 评阅人写三个就多出一行。多出来的先看别的块有没有空着的行：委员只填了四个，
% 六行里空两行，就把这两行匀出去。四块轮着来，一轮各拿一行，拿到够为止；
% 全借完还不够，才把总行数加上去，决议区相应变矮。
%    \begin{macrocode}
\int_new:N \g__hit_res_rev_int    \int_new:N \g__hit_res_chair_int
\int_new:N \g__hit_res_mem_int    \int_new:N \g__hit_res_sec_int
\int_new:N \l__hit_res_need_int
\bool_new:N \l__hit_res_moved_bool
\cs_new_protected:Npn \__hit_res_borrow_one:NN #1#2
  {
    \bool_lazy_and:nnT
      { \int_compare_p:nNn { \l__hit_res_need_int } > { 0 } }
      { \int_compare_p:nNn { #1 } > { \seq_count:N #2 } }
      {
        \int_gdecr:N #1
        \int_decr:N \l__hit_res_need_int
        \bool_set_true:N \l__hit_res_moved_bool
      }
  }
\cs_new_protected:Nn \__hit_res_borrow:
  {
    \int_set:Nn \l__hit_res_need_int
      {
        \int_max:nn { \seq_count:N \g__hit_defense_reviewer_seq  - 2 } { 0 }
        + \int_max:nn { \seq_count:N \g__hit_defense_chair_seq     - 1 } { 0 }
        + \int_max:nn { \seq_count:N \g__hit_defense_member_seq    - 6 } { 0 }
        + \int_max:nn { \seq_count:N \g__hit_defense_secretary_seq - 1 } { 0 }
      }
    \bool_set_true:N \l__hit_res_moved_bool
    \bool_do_while:nn
      {
        \l__hit_res_moved_bool
        && \int_compare_p:nNn { \l__hit_res_need_int } > { 0 }
      }
      {
        \bool_set_false:N \l__hit_res_moved_bool
        \__hit_res_borrow_one:NN \g__hit_res_rev_int
          \g__hit_defense_reviewer_seq
        \__hit_res_borrow_one:NN \g__hit_res_chair_int
          \g__hit_defense_chair_seq
        \__hit_res_borrow_one:NN \g__hit_res_mem_int
          \g__hit_defense_member_seq
        \__hit_res_borrow_one:NN \g__hit_res_sec_int
          \g__hit_defense_secretary_seq
      }
  }
\cs_new_protected:Nn \__hit_resolution_count:
  {
    \int_gset:Nn \g__hit_res_rev_int
      { \int_max:nn { \seq_count:N \g__hit_defense_reviewer_seq } { 2 } }
    \int_gset:Nn \g__hit_res_chair_int
      { \int_max:nn { \seq_count:N \g__hit_defense_chair_seq } { 1 } }
    \int_gset:Nn \g__hit_res_mem_int
      { \int_max:nn { \seq_count:N \g__hit_defense_member_seq } { 6 } }
    \int_gset:Nn \g__hit_res_sec_int
      { \int_max:nn { \seq_count:N \g__hit_defense_secretary_seq } { 1 } }
    \__hit_res_borrow:
  }
%    \end{macrocode}
%
% \subsection{格}
%
% 格里的 \verb|\\| 是换行：词条里用它分行（“评阅人\verb|\\|（根据实际\verb|\\|人数填写）”），
% 行尾另写 \cs{tabularnewline}。表头各格在 \hologo{LuaLaTeX} 下是普通的段，由断行层照 Word 断（“职称（是否
% 博导）”行末的括号压半格，一行排下）；别的引擎 \hologo{TeX} 断行不压标点，这八个字会折成两行，所以按
% \verb|\\| 分行、一行一个行盒，与 Word 一样排成一行。表里的段是 \option{resolution-table}，决议那一格是
% \option{resolution-cell}，两个只差紧排（见上）：范例这张表的段没写行距，
% 跟 \texttt{Normal}，单倍、贴网格，一行一格（iota-hit 的 \texttt{cell-style} 也是这样）；正文的
% \texttt{Normal} 是 1.25 倍，不能直接用。加粗是宋体描边，与 Word 给中易宋体合成粗体同一个做法，
% 见 \file{hit-fonts.dtx} 的 \cs{hit@font@fakebold@zh:}。
%    \begin{macrocode}
\cs_new_protected:Npn \__hit_res_cell:
  { \cs_set_eq:NN \\ \newline }
\cs_new_protected:Npn \__hit_res_size:
  { \bool_if:NTF \g__hit_hass_bool { \wuhao } { \xiaosi } }
\cs_new_protected:Npn \__hit_res_font:
  {
    \docxlike_format_apply_state_else:nn { resolution-table } { \__hit_res_size: }
    \bool_set_eq:NN \l_docxlike_format_par_snap_horizontal_bool
      \l_docxlike_format_snap_horizontal_bool
    \dim_set_eq:NN \l_docxlike_format_par_char_spacing_dim
      \l_docxlike_format_char_spacing_dim
  }
\cs_new_protected:Npn \__hit_res_body_font:
  { \docxlike_format_apply_state_else:nn { resolution-cell } { \__hit_res_size: } }
\cs_new_protected:Nn \__hit_res_bold: { \hit@font@fakebold@zh: }
\cs_new_protected:Npn \__hit_res_head:n #1 { \__hit_res_bold: #1 }
\cs_new_protected:Npn \__hit_res_head_cell:n #1
  {
    \sys_if_engine_luatex:TF
      { \__hit_res_head:n {#1} }
      {
        \seq_set_split:Nnn \l__hit_res_label_seq { \\ } {#1}
        \vbox_top:n
          {
            \offinterlineskip
            \seq_map_function:NN \l__hit_res_label_seq \__hit_res_line:n
          }
      }
  }
%    \end{macrocode}
%
% \emph{每一行先量一遍。} 合并格的标签要在合并出来的高度里居中，而一格里的字长了会折行、那一行就
% 比最小值高，所以先把每一格照它那一栏的宽排进盒里量高，一行取最高的那格与最小值里大的。与 iota-hit
% 的 \texttt{row-of} 一样。量的盒里照表格的规矩排（\cs{docxlike\_table\_begin:}），折行与真排时一致。
%    \begin{macrocode}
\box_new:N \l__hit_res_probe_box
\dim_new:N \l__hit_res_rowh_dim
\cs_new_protected:Npn \__hit_res_measure:nn #1#2
  {
    \vbox_set:Nn \l__hit_res_probe_box
      {
        \__hit_res_font:
        \docxlike_table_begin:
        \box_set_eq:NN \@arstrutbox \strutbox
        \dim_set:Nn \hsize {#1}
        \dim_set_eq:NN \linewidth \hsize
        \centering
        \__hit_res_cell:
        #2 \par
      }
    \dim_set:Nn \l__hit_res_rowh_dim
      {
        \dim_max:nn { \l__hit_res_rowh_dim }
          { \box_ht:N \l__hit_res_probe_box + \box_dp:N \l__hit_res_probe_box }
      }
  }
\cs_new_protected:Npn \__hit_res_record_height:n #1
  {
    \dim_set_eq:NN \l__hit_res_rowh_dim \c__hit_res_row_dim
    \hit@defense@unpack:n {#1}
    \__hit_res_measure:nn { \g__hit_res_wc_dim } { \l__hit_defense_name_tl }
    \__hit_res_measure:nn { \g__hit_res_wd_dim } { \l__hit_defense_title_tl }
    \__hit_res_measure:nn { \g__hit_res_we_dim } { \l__hit_defense_affiliation_tl }
    \__hit_res_measure:nn { \g__hit_res_wf_dim } { \l__hit_defense_discipline_tl }
  }
\cs_new_protected:Npn \__hit_res_head_height:n #1
  {
    \dim_set_eq:NN \l__hit_res_rowh_dim \c__hit_res_row_dim
    \tl_if_blank:nF {#1}
      {
        \__hit_res_measure:nn { \g__hit_res_wb_dim }
          { \__hit_res_head_cell:n { \hit@term@resolution@role } }
      }
    \__hit_res_measure:nn { \g__hit_res_wc_dim } { \__hit_res_head_cell:n { \hit@term@resolution@name } }
    \__hit_res_measure:nn { \g__hit_res_wd_dim } { \__hit_res_head_cell:n { \hit@term@resolution@title } }
    \__hit_res_measure:nn { \g__hit_res_we_dim } { \__hit_res_head_cell:n { \hit@term@resolution@affiliation } }
    \__hit_res_measure:nn { \g__hit_res_wf_dim } { \__hit_res_head_cell:n { \hit@term@resolution@discipline } }
  }
%    \end{macrocode}
%
% 一块的高是它各行的高加上中间的线。填了的行量过，空行就是最小值。
%    \begin{macrocode}
\int_new:N \l__hit_res_i_int
\dim_new:N \l__hit_res_block_dim
\cs_new_protected:Npn \__hit_res_block_height:Nn #1#2
  {
    \dim_zero:N \l__hit_res_block_dim
    \int_zero:N \l__hit_res_i_int
    \seq_map_inline:Nn #1
      {
        \int_incr:N \l__hit_res_i_int
        \__hit_res_record_height:n {##1}
        \dim_add:Nn \l__hit_res_block_dim { \l__hit_res_rowh_dim + \c__hit_res_rule_dim }
      }
    \dim_add:Nn \l__hit_res_block_dim
      {
        ( \c__hit_res_row_dim + \c__hit_res_rule_dim )
        * \int_eval:n { #2 - \l__hit_res_i_int }
        - \c__hit_res_rule_dim
      }
  }
\dim_new:N \g__hit_res_rev_label_dim   \dim_new:N \g__hit_res_com_label_dim
\dim_new:N \g__hit_res_chair_label_dim \dim_new:N \g__hit_res_mem_label_dim
\dim_new:N \g__hit_res_sec_label_dim
\cs_new_protected:Nn \__hit_resolution_heights:
  {
    \__hit_res_head_height:n { }
    \dim_gset_eq:NN \g__hit_res_rev_label_dim \l__hit_res_rowh_dim
    \__hit_res_block_height:Nn \g__hit_defense_reviewer_seq { \g__hit_res_rev_int }
    \dim_gadd:Nn \g__hit_res_rev_label_dim { \c__hit_res_rule_dim + \l__hit_res_block_dim }
    \__hit_res_head_height:n { role }
    \dim_gset_eq:NN \g__hit_res_com_label_dim \l__hit_res_rowh_dim
    \__hit_res_block_height:Nn \g__hit_defense_chair_seq { \g__hit_res_chair_int }
    \dim_gset_eq:NN \g__hit_res_chair_label_dim \l__hit_res_block_dim
    \__hit_res_block_height:Nn \g__hit_defense_member_seq { \g__hit_res_mem_int }
    \dim_gset_eq:NN \g__hit_res_mem_label_dim \l__hit_res_block_dim
    \__hit_res_block_height:Nn \g__hit_defense_secretary_seq { \g__hit_res_sec_int }
    \dim_gset_eq:NN \g__hit_res_sec_label_dim \l__hit_res_block_dim
    \dim_gadd:Nn \g__hit_res_com_label_dim
      {
        \c__hit_res_rule_dim * 3 + \g__hit_res_chair_label_dim
        + \g__hit_res_mem_label_dim + \g__hit_res_sec_label_dim
      }
  }
%    \end{macrocode}
%
% \subsection{合并格}
%
% \env{tabular} 没有纵向合并。合并格的内容放在这一块的头一行里，装成一个挂下去的盒：盒顶对到
% 这一行的顶（基线上方一个撑杆高），深度记零，不撑行高，往下伸进同一栏的空格里；这几行之间那一栏
% 不画横线。盒里先放一条撑杆高的零宽线再退回去，外面那格的高就是撑杆高，与别的格一样。
%    \begin{macrocode}
\box_new:N \l__hit_res_hang_box
\cs_new_protected:Npn \__hit_res_hang:n #1
  {
    \vbox_set_top:Nn \l__hit_res_hang_box
      {
        \hbox:n
          {
            \tex_vrule:D height \box_ht:N \@arstrutbox
              depth \c_zero_dim width \c_zero_dim
          }
        \tex_kern:D - \box_ht:N \@arstrutbox
        \tex_prevdepth:D -1000pt \scan_stop:
        #1
      }
    \box_set_dp:Nn \l__hit_res_hang_box { \c_zero_dim }
    \box_use_drop:N \l__hit_res_hang_box
  }
%    \end{macrocode}
%
% \emph{标签怎么摆照 iota-hit 的 \texttt{\_label}。} 评阅人那一格照 Word 顶端对齐、各行居中，
% 行在词条里用 \verb|\\| 分好（与 Word 在那一栏宽里折的地方一样），一行一个 Word 的行盒，
% 行里的字照 Word 的 run 装（\cs{docxlike\_format\_run\_hbox\_set:Nn}，行头行尾的括号占整格）；
% 别的几块的标签分三种，后两种在合并出来的高度里上下居中：
% \begin{itemize}
% \item 词条里只有一段（“主席”）就排成一行，与别的格一样顶端对齐；英文表的旋转标签居中；
% \item 几个字每字一段（“答辩委员会成员”），一摞排得下就竖排，一字一格，格高 1.04 个字号（范例
%   \texttt{tbRlV} 那一格量的 12.48，iota-hit 取字号加网格余量，默认网格下是同一个数，网格宽的节里
%   不跟着变），字的基线在格顶下一个上升部（中易宋体 220/256\,em），字身的高深不算（\pkg{luatexja} 的
%   JFM 给的 0.88\,em 比上升部高，算进去每字多 0.2）；排不下就横过来连成一行，顶端对齐；
% \item 两个字（“委员”）上下各占一行，中间空几行：这一块排得下几行，减二除五取整，至多两行，
%   两个字加空行排不下就少空。空行是 Word 的行（网格行距），字在行里的位置与别的格一样。
% \end{itemize}
%    \begin{macrocode}
\seq_new:N \l__hit_res_label_seq
\int_new:N \l__hit_res_n_int
\int_new:N \l__hit_res_gap_int
\dim_new:N \l__hit_res_line_dim
\dim_new:N \l__hit_res_pitch_dim
\box_new:N \l__hit_res_text_box
\cs_new_protected:Npn \__hit_res_line:n #1
  {
    \docxlike_format_run_hbox_set:Nn \l__hit_res_text_box { \__hit_res_bold: #1 }
    \hbox_to_wd:nn { \hsize }
      { \hss \box_use:N \@arstrutbox \box_use_drop:N \l__hit_res_text_box \hss }
  }
\cs_new_protected:Npn \__hit_res_top:n #1
  {
    \seq_set_split:Nnn \l__hit_res_label_seq { \\ } {#1}
    \__hit_res_hang:n
      {
        \offinterlineskip
        \seq_map_function:NN \l__hit_res_label_seq \__hit_res_line:n
      }
  }
\cs_new_protected:Npn \__hit_res_center:nn #1#2
  { \__hit_res_hang:n { \vbox_to_ht:nn {#1} { \offinterlineskip \vss #2 \vss } } }
\cs_new_protected:Npn \__hit_res_label:nn #1#2
  {
    \seq_set_split:Nnn \l__hit_res_label_seq { \\ } {#2}
    \int_set:Nn \l__hit_res_n_int { \seq_count:N \l__hit_res_label_seq }
    \dim_set:Nn \l__hit_res_line_dim
      { \box_ht:N \@arstrutbox + \box_dp:N \@arstrutbox }
    \bool_lazy_or:nnTF
      { \int_compare_p:nNn { \l__hit_res_n_int } < { 2 } }
      { \dim_compare_p:nNn { \l__hit_res_line_dim * \l__hit_res_n_int } > {#1} }
      {
        \tl_if_in:nnTF {#2} { \rotatebox }
          {
            \__hit_res_center:nn {#1}
              { \__hit_res_line:n { \seq_use:Nn \l__hit_res_label_seq { } } }
          }
          {
            \__hit_res_hang:n
              { \__hit_res_line:n { \seq_use:Nn \l__hit_res_label_seq { } } }
          }
      }
      {
        \int_compare:nNnTF { \l__hit_res_n_int } > { 2 }
          { \__hit_res_center:nn {#1} { \__hit_res_stack: } }
          { \__hit_res_center:nn {#1} { \__hit_res_pair:n {#1} } }
      }
  }
\tl_const:Nn \c__hit_res_stack_ascent_tl { 0.859375 }
\tl_const:Nn \c__hit_res_stack_pitch_tl { 1.04 }
\cs_new_protected:Npn \__hit_res_stack:
  {
    \dim_set:Nn \l__hit_res_pitch_dim { \c__hit_res_stack_pitch_tl em }
    \seq_map_inline:Nn \l__hit_res_label_seq
      {
        \hbox_set:Nn \l__hit_res_text_box { \__hit_res_bold: ##1 }
        \box_set_ht:Nn \l__hit_res_text_box { \c_zero_dim }
        \box_set_dp:Nn \l__hit_res_text_box { \c_zero_dim }
        \hbox_to_wd:nn { \hsize }
          {
            \hss
            \tex_vrule:D width \c_zero_dim
              height \dim_eval:n { \c__hit_res_stack_ascent_tl em }
              depth \dim_eval:n { \l__hit_res_pitch_dim - \c__hit_res_stack_ascent_tl em }
            \box_use_drop:N \l__hit_res_text_box
            \hss
          }
      }
  }
\cs_new_protected:Npn \__hit_res_pair:n #1
  {
    \int_set:Nn \l__hit_res_gap_int
      {
        \int_min:nn
          {
            \int_max:nn
              {
                (
                  \int_div_truncate:nn
                    { \dim_to_decimal_in_sp:n {#1} }
                    { \dim_to_decimal_in_sp:n { \l__hit_res_line_dim } }
                  - 2
                ) / 5
              }
              { 0 }
          }
          { 2 }
      }
    \bool_while_do:nn
      {
        \int_compare_p:nNn { \l__hit_res_gap_int } > { 0 }
        && \dim_compare_p:nNn
             { \l__hit_res_line_dim * ( \l__hit_res_gap_int + 2 ) } > {#1}
      }
      { \int_decr:N \l__hit_res_gap_int }
    \__hit_res_line:n { \seq_item:Nn \l__hit_res_label_seq { 1 } }
    \prg_replicate:nn { \l__hit_res_gap_int } { \__hit_res_line:n { } }
    \__hit_res_line:n { \seq_item:Nn \l__hit_res_label_seq { 2 } }
  }
%    \end{macrocode}
%
% \subsection{整张表}
%
% 整张表先拼成记号表再放进 \env{tabular}。六栏都是居中的 \texttt{p} 栏，宽是 \texttt{tblGrid} 减两个格
% 边距；评阅人那一块左两栏合成一格（Word 的 \texttt{gridSpan}）。
%    \begin{macrocode}
\tl_new:N \l__hit_res_tbl_tl
\cs_new_protected:Npn \__hit_res_put:n #1
  { \tl_put_right:Nn \l__hit_res_tbl_tl {#1} }
\cs_new_protected:Npn \__hit_res_putx:n #1
  { \tl_put_right:Nx \l__hit_res_tbl_tl {#1} }
\cs_new:Npn \__hit_res_spec:N #1
  {
    > { \exp_not:N \centering \exp_not:N \arraybackslash \exp_not:N \__hit_res_cell: }
    p { \dim_use:N #1 }
  }
%    \end{macrocode}
%
% 一条记录的四格。\pkg{xeCJK} 把汉字设成了活动符，取出来的字段用 \cs{exp\_not:V} 挡住，不展开。
%    \begin{macrocode}
\cs_new_protected:Npn \__hit_res_cells:n #1
  {
    \hit@defense@unpack:n {#1}
    \__hit_res_putx:n
      {
        & \exp_not:V \l__hit_defense_name_tl
        & \exp_not:V \l__hit_defense_title_tl
        & \exp_not:V \l__hit_defense_affiliation_tl
        & \exp_not:V \l__hit_defense_discipline_tl
      }
  }
\cs_new_protected:Npn \__hit_res_head_cells:
  {
    \__hit_res_putx:n
      {
        & \exp_not:N \__hit_res_head_cell:n { \exp_not:o { \hit@term@resolution@name } }
        & \exp_not:N \__hit_res_head_cell:n { \exp_not:o { \hit@term@resolution@title } }
        & \exp_not:N \__hit_res_head_cell:n { \exp_not:o { \hit@term@resolution@affiliation } }
        & \exp_not:N \__hit_res_head_cell:n { \exp_not:o { \hit@term@resolution@discipline } }
        \exp_not:N \tabularnewline
      }
  }
%    \end{macrocode}
%
% 评阅人那一块：表头一行，标签跨下来；每行左边是合成一格的空格。
%    \begin{macrocode}
\cs_new_protected:Nn \__hit_res_reviewer_rows:
  {
    \__hit_res_putx:n
      {
        \exp_not:N \multicolumn { 2 } { | \__hit_res_spec:N \g__hit_res_wab_dim | }
          {
            \exp_not:N \__hit_res_top:n { \exp_not:o { \hit@term@resolution@reviewer } }
          }
      }
    \__hit_res_head_cells:
    \__hit_res_hrule:nN { 3 } \c_true_bool
    \int_zero:N \l__hit_res_i_int
    \seq_map_inline:Nn \g__hit_defense_reviewer_seq
      {
        \int_incr:N \l__hit_res_i_int
        \__hit_res_reviewer_empty:
        \__hit_res_cells:n {##1}
        \__hit_res_row_end:nnnN { \g__hit_res_rev_int } { \__hit_res_hrule:nN { 1 } \c_false_bool } { 3 }
          \c_true_bool
      }
    \int_step_inline:nn { \g__hit_res_rev_int - \l__hit_res_i_int }
      {
        \int_incr:N \l__hit_res_i_int
        \__hit_res_reviewer_empty:
        \__hit_res_put:n { & & & & }
        \__hit_res_row_end:nnnN { \g__hit_res_rev_int } { \__hit_res_hrule:nN { 1 } \c_false_bool } { 3 }
          \c_true_bool
      }
  }
\cs_new_protected:Npn \__hit_res_reviewer_empty:
  {
    \__hit_res_putx:n
      { \exp_not:N \multicolumn { 2 } { | \__hit_res_spec:N \g__hit_res_wab_dim | } { } }
  }
%    \end{macrocode}
%
% \emph{横线都占高。} Word 的每道横框线占它自己的粗细，一块里合并格那几栏不画线，别的栏照画，
% 这一行的高照样多半磅。\cs{cline} 画完往回退一道线宽，下一行压在线上，行距就少了半磅；而补回那
% 半磅的话，没画线的那几栏竖线会断一截。所以不用 \cs{cline}，另起一行：每格 \cs{omit}，要画的栏
% 放一道横线的 \cs{leaders}，格界放 \pkg{docxlike} 的竖线，行高就是线宽，竖线接得上。
% \meta{栏} 是从第几栏起画，\meta{合并} 是左两栏合成了一格（评阅人那一块，中间没有竖线）。
%
% 整道的横线（表顶、评阅人那一块下面、表底）也这样画，不用 \cs{hline}：外框的竖线骑在表的外沿上，
% 一半在表外，\cs{hline} 只画到表的外沿，每道横线与外框相交处就缺一个线宽一半的角，决议框接上来的
% 地方看得出一道凹缝。这样画的那一行外沿也有竖线，角是满的。
%    \begin{macrocode}
\cs_new:Npn \__hit_res_leaders:
  {
    \exp_not:N \leaders \exp_not:N \hrule height \dim_use:N \c__hit_res_rule_dim
    \exp_not:N \hfill
  }
\cs_new_protected:Npn \__hit_res_hrule:nN #1#2
  {
    \__hit_res_putx:n
      {
        \exp_not:N \omit \exp_not:N \vline
          \int_compare:nNnT {#1} < { 2 } { \__hit_res_leaders: }
        & \exp_not:N \omit
          \bool_if:NF #2 { \exp_not:N \vline }
          \int_compare:nNnT {#1} < { 3 } { \__hit_res_leaders: }
        & \exp_not:N \omit \exp_not:N \vline \__hit_res_leaders:
        & \exp_not:N \omit \exp_not:N \vline \__hit_res_leaders:
        & \exp_not:N \omit \exp_not:N \vline \__hit_res_leaders:
        & \exp_not:N \omit \exp_not:N \vline \__hit_res_leaders: \exp_not:N \vline
        \exp_not:N \tabularnewline
      }
  }
%    \end{macrocode}
%
% 行尾：一块的最后一行下面照 \meta{块末} 画（整道横线或下一块的那一截），别的行只画 \meta{栏} 到第六栏。
%    \begin{macrocode}
\cs_new_protected:Npn \__hit_res_row_end:nnnN #1#2#3#4
  {
    \__hit_res_put:n { \tabularnewline }
    \int_compare:nNnTF { \l__hit_res_i_int } = {#1}
      {#2}
      { \__hit_res_hrule:nN {#3} #4 }
  }
%    \end{macrocode}
%
% 委员会那一块：表头一行，委员会的标签跨表头与下面三块；主席、委员、秘书各一块，职务标签跨这一块。
%    \begin{macrocode}
\cs_new_protected:Npn \__hit_res_committee_rows:NnNnn #1#2#3#4#5
  {
    \int_zero:N \l__hit_res_i_int
    \seq_map_inline:Nn #1
      {
        \int_incr:N \l__hit_res_i_int
        \__hit_res_role:nNn {#2} #3 {#4}
        \__hit_res_cells:n {##1}
        \__hit_res_row_end:nnnN {#2} {#5} { 3 } \c_false_bool
      }
    \int_step_inline:nn { #2 - \l__hit_res_i_int }
      {
        \int_incr:N \l__hit_res_i_int
        \__hit_res_role:nNn {#2} #3 {#4}
        \__hit_res_put:n { & & & & }
        \__hit_res_row_end:nnnN {#2} {#5} { 3 } \c_false_bool
      }
  }
\cs_new_protected:Npn \__hit_res_role:nNn #1#2#3
  {
    \__hit_res_put:n { & }
    \int_compare:nNnT { \l__hit_res_i_int } = { 1 }
      {
        \__hit_res_putx:n
          {
            \exp_not:N \__hit_res_label:nn { \dim_use:N #2 } { \exp_not:o {#3} }
          }
      }
  }
\cs_new_protected:Nn \__hit_res_build:
  {
    \tl_clear:N \l__hit_res_tbl_tl
    \__hit_res_hrule:nN { 1 } \c_false_bool
    \__hit_res_reviewer_rows:
    \__hit_res_putx:n
      {
        \exp_not:N \__hit_res_label:nn
          { \dim_use:N \g__hit_res_com_label_dim }
          { \exp_not:o { \hit@term@resolution@committee } }
        & \exp_not:N \__hit_res_head_cell:n { \exp_not:o { \hit@term@resolution@role } }
      }
    \__hit_res_head_cells:
    \__hit_res_hrule:nN { 2 } \c_false_bool
    \__hit_res_committee_rows:NnNnn \g__hit_defense_chair_seq { \g__hit_res_chair_int }
      \g__hit_res_chair_label_dim { \hit@term@resolution@chair }
      { \__hit_res_hrule:nN { 2 } \c_false_bool }
    \__hit_res_committee_rows:NnNnn \g__hit_defense_member_seq { \g__hit_res_mem_int }
      \g__hit_res_mem_label_dim { \hit@term@resolution@member }
      { \__hit_res_hrule:nN { 2 } \c_false_bool }
    \__hit_res_committee_rows:NnNnn \g__hit_defense_secretary_seq { \g__hit_res_sec_int }
      \g__hit_res_sec_label_dim { \hit@term@resolution@secretary }
      { \__hit_res_hrule:nN { 1 } \c_false_bool }
  }
%    \end{macrocode}
%
% 表放在一只与版心同宽的盒里两头 \cs{hss}，比版心宽出的那一截两边探出。表顶接着上一段排，不加行间胶：
% Word 的表紧跟着前一段的段后。
%    \begin{macrocode}
\cs_new_protected:Nn \__hit_res_table:
  {
    \group_begin:
      \dim_set:Nn \arrayrulewidth { \c__hit_res_rule_dim }
      \dim_set_eq:NN \tabcolsep \c__hit_res_inset_dim
      \__hit_res_font:
      \docxlikesetup
        { table-properties = { row = { specify-height = \dim_use:N \c__hit_res_row_dim } } }
      \nointerlineskip
      \hbox_to_wd:nn { \linewidth }
        {
          \hss
          \use:x
            {
              \exp_not:N \begin { tabular } [ t ]
                {
                  | \__hit_res_spec:N \g__hit_res_wa_dim
                  | \__hit_res_spec:N \g__hit_res_wb_dim
                  | \__hit_res_spec:N \g__hit_res_wc_dim
                  | \__hit_res_spec:N \g__hit_res_wd_dim
                  | \__hit_res_spec:N \g__hit_res_we_dim
                  | \__hit_res_spec:N \g__hit_res_wf_dim |
                }
            }
          \tl_use:N \l__hit_res_tbl_tl
          \end { tabular }
          \hss
        }
    \group_end:
  }
%    \end{macrocode}
%
% \subsection{决议区}
%
% 范例里是这张表的最后一行，六栏合成一格，行高最小值 2835 缇，格里是标签一段加一串空段。
% Word 的行能跨页，\env{tabular} 的不能，所以这一格在表下另画：左右两道竖线接着表的外框，底下一道横线，
% 顶线就是表的最后一道横线；跨到下一页的那一截自己补上顶线。
%
% 高度照 iota-hit：排到这一页剩下的整数个网格行；正文比这高就照正文排，到页底切开，下一页接着排，
% 最后一截照正文的高。一行都不剩就整格挪到下一页，在那一页照样排满。
%
% 格里：标签顶格加粗，正文每段首行缩进两个字（与 iota-hit 一样），格边距与表的一样，字从版心左缘起。
%    \begin{macrocode}
\box_new:N \l__hit_res_body_box
\box_new:N \l__hit_res_part_box
\bool_new:N \l__hit_res_first_bool
\bool_new:N \l__hit_res_fill_bool
\dim_new:N \l__hit_res_rest_dim
\dim_new:N \l__hit_res_lead_dim
\cs_new_protected:Nn \__hit_res_lead:
  {
    \dim_compare:nNnTF { \g_docxlike_page_lead_dim } > { \c_zero_dim }
      { \dim_set_eq:NN \l__hit_res_lead_dim \g_docxlike_page_lead_dim }
      { \dim_set_eq:NN \l__hit_res_lead_dim \baselineskip }
  }
\cs_new_protected:Npn \__hit_res_lines:n #1
  {
    \dim_set:Nn \l__hit_res_rest_dim
      {
        \fp_to_dim:n
          {
            floor ( \dim_to_fp:n {#1} / \dim_to_fp:n { \l__hit_res_lead_dim } )
            * \dim_to_fp:n { \l__hit_res_lead_dim }
          }
      }
    \dim_compare:nNnT { \l__hit_res_rest_dim } < { \c_zero_dim }
      { \dim_zero:N \l__hit_res_rest_dim }
  }
\cs_new_protected:Nn \__hit_res_body:
  {
    \vbox_set:Nn \l__hit_res_body_box
      {
        \__hit_res_body_font:
        \docxlike_table_begin:
        \dim_set_eq:NN \hsize \g__hit_res_text_dim
        \dim_set_eq:NN \linewidth \hsize
        \skip_zero:N \leftskip
        \skip_zero:N \rightskip
        \skip_set:Nn \parfillskip { 0pt plus 1fil }
        \dim_zero:N \parindent
        { \__hit_res_bold: \hit@term@resolution@text }
        \par
        \dim_compare:nNnTF { \g_docxlike_page_char_dim } > { \c_zero_dim }
          { \dim_set:Nn \parindent { 2 \g_docxlike_page_char_dim } }
          { \dim_set:Nn \parindent { 2em } }
        \g__hit_defense_resolution_tl
        \par
      }
  }
%    \end{macrocode}
%
% 画框。竖线骑在表的外框线上（格边距里扣掉半道线），框里的盒给了高度就顶端对齐排到那么高。
%    \begin{macrocode}
\cs_new_protected:Npn \__hit_res_frame:nN #1#2
  {
    \nointerlineskip
    \hbox_to_wd:nn { \linewidth }
      {
        \hss
        \vbox:n
          {
            \bool_if:NF \l__hit_res_first_bool
              { \tex_hrule:D height \c__hit_res_rule_dim }
            \hbox:n
              {
                \tex_vrule:D width \c__hit_res_rule_dim
                \skip_horizontal:n { \tabcolsep - \c__hit_res_rule_dim / 2 }
                \tl_if_blank:nTF {#1}
                  { \box_use:N #2 }
                  { \vbox_to_ht:nn {#1} { \box_use:N #2 \vss } }
                \skip_horizontal:n { \tabcolsep - \c__hit_res_rule_dim / 2 }
                \tex_vrule:D width \c__hit_res_rule_dim
              }
            \tex_hrule:D height \c__hit_res_rule_dim
          }
        \hss
      }
  }
%    \end{macrocode}
%
% 这一页剩多少要等表进了主竖直列表才读得到：先发一个 10000 的罚分让分页器跑一趟，再拿
% \cs{pagegoal} 减 \cs{pagetotal}，扣掉框的底线，向下取整到网格行。
%    \begin{macrocode}
\cs_new_protected:Nn \__hit_res_page_rest:
  {
    \tex_penalty:D 10000 \scan_stop:
    \__hit_res_lines:n
      { \pagegoal - \pagetotal - \pagedepth - \c__hit_res_rule_dim }
  }
\cs_new_protected:Npn \__hit_res_conclusion_page:n #1
  {
    \dim_compare:nNnTF
      { \box_ht:N \l__hit_res_body_box + \box_dp:N \l__hit_res_body_box } > {#1}
      {
        \vbox_set_split_to_ht:NNn \l__hit_res_part_box \l__hit_res_body_box {#1}
        \__hit_res_frame:nN {#1} \l__hit_res_part_box
        \newpage
        \bool_set_false:N \l__hit_res_first_bool
        \bool_set_false:N \l__hit_res_fill_bool
        \__hit_res_lines:n { \textheight - 2 \c__hit_res_rule_dim }
        \exp_args:NV \__hit_res_conclusion_page:n \l__hit_res_rest_dim
      }
      {
        \bool_if:NTF \l__hit_res_fill_bool
          { \__hit_res_frame:nN {#1} \l__hit_res_body_box }
          { \__hit_res_frame:nN { } \l__hit_res_body_box }
      }
  }
\cs_generate_variant:Nn \__hit_res_conclusion_page:n { V }
\cs_new_protected:Nn \__hit_resolution_sheet:
  {
    \__hit_resolution_calc:
    \__hit_resolution_count:
    \__hit_resolution_heights:
    \__hit_res_build:
    \par
    \__hit_res_table:
    \__hit_res_lead:
    \__hit_res_body:
    \skip_zero:N \splittopskip
    \bool_set_true:N \l__hit_res_first_bool
    \bool_set_true:N \l__hit_res_fill_bool
    \__hit_res_page_rest:
    \dim_compare:nNnT { \l__hit_res_rest_dim } < { \l__hit_res_lead_dim }
      {
        \newpage
        \bool_set_false:N \l__hit_res_first_bool
        \__hit_res_lines:n { \textheight - 2 \c__hit_res_rule_dim }
      }
    \__hit_res_conclusion_page:V \l__hit_res_rest_dim
  }
%    \end{macrocode}
%
% \begin{macro}{\hit@resolution@place}
% 三种用法：给 PDF 就贴那个 PDF，带星号连页码一起排；不给就按
% \texttt{defense} 族的数据排出来，没数据就是一张空表。
%
% 学位这一层的判断在部件那一层（\cs{defense} 的 \option{shown}，见 \file{src/engine/hit-parts.dtx}），
% 所以这里不再自己查 \option{doctor}。
%
% \textbf{为什么要拆成两个。} \cs{resolution} 是 v3.1e 的写法，v3.2a 起叫
% \cs{defense}，用户再调旧名要给一句改名提示。
% 但类自己排这一页走的也是同一段代码，提示不能跟着响。所以排版的活留在这个
% 内部命令里，\cs{resolution} 只负责提示再转发过来。
%    \begin{macrocode}
\NewDocumentCommand \hit@resolution@place { s o }
  {
    \hit@clear@page
    \IfNoValueTF {#2}
      {
        \hit@appendix@chapter * { \hit@term@resolution@heading@zh }
          [ \hit@term@resolution@heading@en ]
        \__hit_resolution_sheet:
      }
      {
        \phantomsection
        \addcontentsline { toc } { chapter } { \hit@term@resolution@heading@zh }
        \addcontentsline { toe } { chapter }
          {
            \texorpdfstring
              { \bfseries \hit@term@resolution@heading@en } { \hit@term@resolution@heading@en }
          }
        \IfBooleanTF {#1}
          {
            \includepdf
              [ fitpaper=true , pages=- ,
                pagecommand={\thispagestyle{hit@scanpage}} ] {#2}
          }
          {
            \includepdf
              [ fitpaper=true , pages=- ,
                pagecommand={\thispagestyle{empty}} ] {#2}
          }
      }
  }
%    \end{macrocode}
% \end{macro}
%
% \begin{macro}{\defense,\resolution}
% \cs{defense}\oarg{选项} 与 iota-hit 的 \texttt{defense()} 对应：\option{shown}、贴扫描件的 \option{file} 与
% \option{page-number}。照 iota-hit，终稿里除博士后都排（\texttt{auto}）；原先只有博士排。名单与决议正文
% 也照 iota-hit 是这一页的参数：\option{reviewers}、\option{members} 是一串记录，\option{chair}、\option{secretary}
% 是一条记录，\option{resolution} 是决议正文。一条记录照旧四个字段：\option{name}、\option{title}、
% \option{affiliation}、\option{discipline}。v3.1e 的 \cs{resolution}\oarg{文件}（星号排页码）照旧能用，转发过来并提示改名。
%    \begin{macrocode}
\keys_define:nn { hit / part / defense }
  {
    shown .code:n = { \hit@part@set@shown:n {#1} } ,
    file .tl_set:N = \l_hit_part_file_tl ,
    page-number .bool_set:N = \l_hit_part_page_number_bool ,
    page-number .default:n = true ,
    reviewers .code:n =
      { \clist_map_inline:nn {#1} { \seq_gput_right:Nn \g__hit_defense_reviewer_seq {##1} } } ,
    chair .code:n = { \seq_gput_right:Nn \g__hit_defense_chair_seq {#1} } ,
    members .code:n =
      { \clist_map_inline:nn {#1} { \seq_gput_right:Nn \g__hit_defense_member_seq {##1} } } ,
    secretary .code:n = { \seq_gput_right:Nn \g__hit_defense_secretary_seq {#1} } ,
    resolution .code:n = { \tl_gset:Nn \g__hit_defense_resolution_tl {#1} } ,
  }
\NewDocumentCommand \defense { +O{} }
  {
    \tl_clear:N \l_hit_part_file_tl
    \bool_set_false:N \l_hit_part_page_number_bool
    \hit@part@opts:nn { defense } {#1}
    \hit@part:nVn { defense } \l_hit_part_shown_tl
      {
        \hit@matter@back@enter:
        \hit@part@scanned:N \hit@resolution@place
      }
  }
\NewDocumentCommand \resolution { s o }
  {
    \hit@warn@renamed@to:nn { resolution } { defense }
    \use:e
      {
        \exp_not:N \defense
          [ \IfValueT {#2} { file = { \exp_not:n {#2} } , } \IfBooleanT {#1} { page-number } ]
      }
  }
%    \end{macrocode}
% \end{macro}
%
%    \begin{macrocode}
  }
%</final-resolution>
%    \end{macrocode}
%
% \Finale
